% element: 10-s064:e05
\BeleskeID{10-s064}
Početni uslov u punom zapisu je $v_O(t_0^+)=v_C(t_0^+)=v_C(t_0^-)=V_L$. U zasićenom ekvivalentu kondenzator vidi paralelnu vezu otpornika.

Kontura daje
\[v_{RB}(t)+V_{BES}-V_{CES}-v_{RC}(t)=0,\]
\[R_Bi_{RB}(t)+V_{BES}-V_{CES}-R_Ci_{RC}(t)=0.\]
U asimptoti kondenzator ne vodi, pa $i_{RB}(\infty)+i_{RC}(\infty)=0$. Iz tih uslova dobijamo asimptotsku struju kolektorskog otpornika. Izlazni napon potom računamo kao
\[v_O(\infty)=V_{CC}-R_Ci_{RC}(\infty)-V_{CES}-V_D.\]
Modelna asimptota ne znači da će ceo prelaz ostati u zasićenju. U određenom trenutku tranzistor prelazi u aktivni režim; od tog trenutka koristi se aktivni model sa novom početnom vrednošću.
